\section{Exact values and optimal positional strategies}
\label{sec:full-solution}

We now use the winning-set algorithm to recover the numerical values
and strategies of the game. For the input model of
Theorem~\ref{thm:main}, write
\[
 \operatorname{val}_w(v)
 =\sup_\sigma\inf_\tau \MP_w(\pi_{v,\sigma,\tau}),
\]
where the strategies may depend on the full finite history.

\begin{corollary}[Full game solution]\label{cor:full-solution}
A uniform deterministic algorithm returns the exact value
$\operatorname{val}_w(v)$ as a reduced rational for every $v\in V$
and one positional strategy $\sigma^\star$ for Max and one positional
strategy $\tau^\star$ for Min such that, simultaneously for every
$v\in V$ and every pair of history-dependent strategies $\sigma,\tau$,
\[
 \MP_w(\pi_{v,\sigma^\star,\tau})\ge\operatorname{val}_w(v),
 \qquad
 \MP_w(\pi_{v,\sigma,\tau^\star})\le\operatorname{val}_w(v).
\]
The complete computation takes
$2^{O((\log_2(L+2))^2)}$ bit operations.
\end{corollary}

\begin{proof}
\noindent\textbf{Positional values.}\enspace
The uniform form of Ehrenfeucht--Mycielski positional determinacy
\cite{EhrenfeuchtMycielski1979}, stated explicitly in
\cite[Theorem~1 and the following paragraph]{BrimEtAl2011}, provides
one positional strategy for each player attaining the value from all
vertices at once, against arbitrary history-dependent opponents.
The latter source uses an edge relation without parallel entries.
For our convention, replace each explicit edge $e:u\to v$ by
$u\to x_e\to v$ through a fresh vertex with a single outgoing edge,
with weights $2w(e)$ and $0$. Keep ownership at the original vertices;
the owner of $x_e$ is immaterial. The resulting graph has no parallel
edges or loops. Expanded histories retain the identity of every
original edge. Their averages at even prefixes equal the original
averages, and the odd-prefix discrepancy tends to zero. Thus payoffs
and the uniform positional guarantees transfer back to our graph,
including loops and parallel entries.

Set $W_{\mathrm{in}}=\max_{e\in E}|w(e)|$. Fix both uniform optimal
positional strategies. From any vertex the resulting play eventually
repeats a simple directed cycle of at most $n$ edges in the original
graph, and its mean is the value at that vertex. Consequently every
value has a reduced representation $a/b$ with
\[
 1\le b\le n,\qquad |a|\le bW_{\mathrm{in}},
 \qquad -W_{\mathrm{in}}\le a/b\le W_{\mathrm{in}}.
\]
The same bounds hold after restricting outgoing choices, provided at
least one outgoing edge remains at every vertex.

\par\medskip\noindent\textbf{Threshold queries and exact reconstruction.}\enspace
We use the classical reduction from values to threshold queries
\cite[Section~5.2]{BrimEtAl2011}, keeping track of binary lengths.
For integers $p$ and $q>0$, replace each weight by
$w_{p/q}(e)=qw(e)-p$. For every infinite play,
\[
 \MP_{w_{p/q}}(\pi)=q\MP_w(\pi)-p.
\]
Positional attainment therefore makes the winning set returned by
Theorem~\ref{thm:main} exactly
$\{v:\operatorname{val}_w(v)\ge p/q\}$.

Choose
$K=1+\lceil\log_2(2(W_{\mathrm{in}}+1)n^2)\rceil$, computable from
integer bit lengths. For each vertex, start with the closed interval
$[-W_{\mathrm{in}},W_{\mathrm{in}}]$ and make $K$ bisections. Query
its midpoint and move the lower endpoint to the midpoint if the
vertex belongs to the returned set; otherwise move the upper endpoint.
The final interval contains its value and has width
$2W_{\mathrm{in}}2^{-K}<1/n^2$. Distinct reduced rationals with
positive denominators at most $n$ differ by at least $1/n^2$.
To reconstruct the value from an interval $[\ell,h]$, for each
$1\le b\le n$ test the integer $a=\lceil b\ell\rceil$ for
$a/b\le h$, reduce the retained fractions, and remove repetitions.
Since $b(h-\ell)<1$, this tests every possible numerator.
The value bound above guarantees exactly one resulting rational.
Call this value-vector procedure $\mathsf{Values}(H)$ for any
restricted game $H$; the original $n,W_{\mathrm{in}},K$ remain valid
for all such calls.

\par\medskip\noindent\textbf{Separate strategy self-reductions.}\enspace
An edge joining vertices of the same value need not be optimal.
For example, let a Max vertex $u$ have a zero-weight self-loop and a
zero-weight edge to a vertex $v$ whose only edge is a self-loop of
weight $1$. Both vertices have value $1$, but staying at $u$ yields
payoff $0$. We therefore test restrictions of the current game and
preserve its entire value vector.

First compute $\mathbf v=\mathsf{Values}(G)$. For Max, start from a
fresh copy $H=G$ and process the Max vertices in any fixed order.
At a vertex $u$, form a trial $H_e$ for each outgoing edge $e$ by
keeping only $e$ at $u$. Keep a trial for which
$\mathsf{Values}(H_e)=\mathbf v$, and continue from that trial.
Such an edge always exists. Indeed, take the edge chosen at $u$ by a
globally optimal positional Max strategy of the current $H$.
Restricting Max's choices cannot increase any value, while that
strategy remains available and secures the current entire value
vector, so no value decreases either. This proves the invariant
$\mathsf{Values}(H)=\mathbf v$.

At the end, Max has one allowed edge at every owned vertex, whereas
all of Min's original choices remain. The resulting single Max
strategy therefore secures $\mathbf v$ from every vertex against
every Min strategy. To recover Min's strategy, start again from a
fresh copy of the original $G$, not from the Max-restricted game.
At each Min vertex keep an outgoing edge whose trial preserves
$\mathbf v$. A globally optimal positional Min strategy of the
current game supplies such an edge: restricting Min cannot decrease
values, and that strategy prevents any increase. The final Min
strategy is optimal against every original Max strategy. Thus the two
separately recovered strategies have the simultaneous guarantees in
the statement.

\par\medskip\noindent\textbf{Call and bit bounds.}\enspace
Put $m=|E|$. At a vertex, at most its original outdegree many trials
are tested, so the two self-reductions together use at most $m$ trial
value-vector computations. Including the original computation, there
are at most
\begin{equation}
 Q=nK(1+m)
 \label{eq:threshold-call-bound}
\end{equation}
threshold calls. Every midpoint has denominator dividing $2^K$ and
a numerator of $O(\log(W_{\mathrm{in}}+1)+K)$ bits. Each transformed
weight has the same order of bit length. Since
$K=O(L+\log(n+1))$, the full encoding length of every queried game
is $O(L+mK)$, polynomial in $L+1$. Reconstruction, reduced-rational
comparison, and copying restricted edge lists also have polynomial
bit cost. The explicit encoding bounds $n,m$ polynomially in $L+1$,
so $Q$ is polynomial. Applying Theorem~\ref{thm:main} to these
polynomial-length games and multiplying by $Q$ gives
$2^{O((\log_2(L+2))^2)}$ bit operations, including the polynomial-size
value and strategy output.
\end{proof}

The preceding reduction also applies to the companion algorithm.
Here the additional issue is that threshold queries depend on earlier
answers, so the winning-set error must be controlled at each query.

\begin{corollary}[Randomized full game solution]
\label{cor:randomized-full-solution}
For the game model of Theorem~\ref{thm:main}, a uniform randomized
algorithm returns all exact values and both globally optimal positional
strategies specified in Corollary~\ref{cor:full-solution} with
probability at least $7/8$ for the entire output. Its running time is
$2^{O((\log_2(L+2))^2)}$ bit operations on every random tape.
\end{corollary}

\begin{proof}
The companion's Theorem~1.1 \cite{OpenAI2026Randomized} returns the
whole winning set correctly with probability at least $7/8$ and has
the quasipolynomial bit bound on every tape. Use the reduction above,
whose number of threshold queries is bounded by the polynomial $Q$
in Equation~\eqref{eq:threshold-call-bound}. For each query, take
$r=2\lceil\log_2(8Q)\rceil+1$ independent fresh runs and use the
coordinatewise majority of their winning sets. Conditional on the
previous history, the current query is fixed. The amplified set is
correct whenever more than half of the runs return the whole correct
set, and its conditional failure probability is at most
\[
 \sum_{j=\lceil r/2\rceil}^{r}\binom rj(1/8)^j
 \le 2^r(1/8)^{r/2}=2^{-r/2}<\frac1{8Q}.
\]
A first-error union bound over the adaptive queries makes every
threshold answer correct with probability at least $7/8$. On this
event the preceding deterministic argument gives all exact values and
both globally optimal strategies. On other tapes, if reconstruction
does not yield a unique rational or no value-preserving trial is
found, the reduction stops with zero at every value coordinate and
the first listed outgoing edge at every owned vertex.
All loops have the stated caps even then. Since $r=O(\log(L+2))$ and
all queried encodings retain their polynomial bound on every tape,
the total bit bound remains quasipolynomial on every tape, including
random-bit generation. This transfer uses only the winning-set
interface.
\end{proof}
